\input{../../../templates/es/preambulo.tex}

% Los listados se toman de src/ con \lstinputlisting; la diapositiva es el archivo.
\lstset{
    basicstyle=\ttfamily\fontsize{8pt}{9.6pt}\selectfont,
    breaklines=true,
    breakatwhitespace=true,
    columns=fullflexible,
    keepspaces=true,
    xleftmargin=0.3em,
    framesep=2pt,
    aboveskip=0pt,
    belowskip=0pt
}

\newcommand{\marcoresultado}[3]{%
    \begin{frame}{#1}
        \begin{center}
            \includegraphics[height=0.62\textheight,width=\textwidth,keepaspectratio]{#2}
        \end{center}
        \vspace{0.25em}
        {\small #3\par}
    \end{frame}%
}

\newcommand{\marcoafirmacion}[2]{%
    \begin{frame}{#1}
        {\small #2\par}
    \end{frame}%
}

\date{}

\begin{document}

\tituloleccion{09}{Optimización Combinatoria Binaria}

\section{Esta lección}

\begin{frame}{Esta lección}
La Lección 08 mezcló tipos en una caja que no se podía dibujar. Hoy el cromosoma es bits, y tres problemas comparten los operadores.

\vspace{0.6em}
\begin{itemize}
    \item Mochila, horarios, radar: los mismos bits, distinto significado, distinta factibilidad.
    \item Reparar cambia la distribución que se busca; no es plomería gratis.
    \item Las instancias chicas hay que comprobarlas exactamente.
\end{itemize}
\end{frame}

% ================= knapsack_01_la_instancia =================
\begin{frame}{Tres modelos binarios}
\small
\textbf{Mochila:}\quad
\(\max\sum_i v_ix_i\) sujeto a \(\sum_i w_ix_i\le C\).\par
\medskip
\textbf{Radar:}\quad
\(\min\;U(x)(|S|+1)+\sum_i x_i\). Como a lo sumo hay \(|S|\) radares,
reducir en uno el número de objetivos sin cobertura domina cualquier cambio de cantidad.\par
\medskip
\textbf{Horarios:}\quad
\(\min\sum_{eds}c_{eds}x_{eds}\), con
\(\sum_e x_{eds}=d_s\) y \(\sum_{ds}x_{eds}\le5\).
Reparar cambia la distribución que realmente se busca.
\end{frame}

\section{La instancia}

\begin{frame}{La instancia}
Cada bit responde a una pregunta: tomar este artículo o dejarlo. Antes de introducir un
AG, enumere esta instancia deliberadamente pequeña para que las ejecuciones posteriores tengan una
respuesta honesta contra la cual comparar.
\end{frame}

\marcoafirmacion{La instancia}{%
    La instancia de 12 artículos tiene 4,096 cromosomas; el óptimo exacto tiene un valor de 50 con un peso de 20}

% ================= knapsack_02_poblacion_aleatoria =================
\section{Una población aleatoria}

\begin{frame}{Una población aleatoria}
Una población aleatoria
\end{frame}

\begin{frame}[fragile]{(1) Individuo}
\lstinputlisting[basicstyle=\ttfamily\fontsize{8pt}{9.6pt}\selectfont,firstline=59,lastline=69]{../src/knapsack_02_poblacion_aleatoria.py}
\end{frame}

\begin{frame}[fragile]{(2) individuo\_aleatorio()}
\lstinputlisting[basicstyle=\ttfamily\fontsize{8pt}{9.6pt}\selectfont,firstline=72,lastline=77]{../src/knapsack_02_poblacion_aleatoria.py}
\end{frame}

\begin{frame}[fragile]{(3) TAMANO\_POBLACION}
\lstinputlisting[basicstyle=\ttfamily\fontsize{8pt}{9.6pt}\selectfont,firstline=80,lastline=82]{../src/knapsack_02_poblacion_aleatoria.py}
\end{frame}

\begin{frame}[fragile]{(4) el censo}
\lstinputlisting[basicstyle=\ttfamily\fontsize{8pt}{9.6pt}\selectfont,firstline=85,lastline=90]{../src/knapsack_02_poblacion_aleatoria.py}
\end{frame}

\marcoafirmacion{Una población aleatoria}{%
    233 de 1,000 cromosomas aleatorios son factibles; el mejor valor aleatorio es 47}

% ================= knapsack_03_reparacion =================
\section{La reparación crea presión de selección}

\begin{frame}{La reparación crea presión de selección}
La reparación crea presión de selección
\end{frame}

\begin{frame}[fragile]{(1) reparar()}
\lstinputlisting[basicstyle=\ttfamily\fontsize{8pt}{9.6pt}\selectfont,firstline=77,lastline=88]{../src/knapsack_03_reparacion.py}
\end{frame}

\begin{frame}[fragile]{(2) la comparación}
\lstinputlisting[basicstyle=\ttfamily\fontsize{8pt}{9.6pt}\selectfont,firstline=91,lastline=97]{../src/knapsack_03_reparacion.py}
\end{frame}

\marcoafirmacion{La reparación crea presión de selección}{%
    La reparación hace factibles a 1,000/1,000 y cambia 767 cromosomas}

% ================= knapsack_04_seleccion_y_cruza =================
\section{Selección y cruza}

\begin{frame}{Selección y cruza}
Selección y cruza
\end{frame}

\begin{frame}[fragile]{(1) seleccion\_torneo()}
\lstinputlisting[basicstyle=\ttfamily\fontsize{8pt}{9.6pt}\selectfont,firstline=85,lastline=88]{../src/knapsack_04_seleccion_y_cruza.py}
\end{frame}

\begin{frame}[fragile]{(2) cruza()}
\lstinputlisting[basicstyle=\ttfamily\fontsize{8pt}{9.6pt}\selectfont,firstline=91,lastline=97]{../src/knapsack_04_seleccion_y_cruza.py}
\end{frame}

\begin{frame}[fragile]{(3) una\_generacion()}
\lstinputlisting[basicstyle=\ttfamily\fontsize{8pt}{9.6pt}\selectfont,firstline=100,lastline=109]{../src/knapsack_04_seleccion_y_cruza.py}
\end{frame}

\marcoafirmacion{Selección y cruza}{%
    Una generación eleva el valor medio de 39.00 a 40.73; todos los hijos almacenados siguen siendo factibles}

% ================= knapsack_05_la_busqueda_completa =================
\section{La búsqueda completa}

\begin{frame}{La búsqueda completa}
La búsqueda completa
\end{frame}

\begin{frame}[fragile]{(1) mutar()}
\lstinputlisting[basicstyle=\ttfamily\fontsize{8pt}{9.6pt}\selectfont,firstline=97,lastline=102]{../src/knapsack_05_la_busqueda_completa.py}
\end{frame}

\begin{frame}[fragile]{(2) ejecutar()}
\lstinputlisting[basicstyle=\ttfamily\fontsize{8pt}{9.6pt}\selectfont,firstline=113,lastline=127]{../src/knapsack_05_la_busqueda_completa.py}
\end{frame}

\begin{frame}[fragile]{(3) el reporte de brecha}
\lstinputlisting[basicstyle=\ttfamily\fontsize{8pt}{9.6pt}\selectfont,firstline=130,lastline=135]{../src/knapsack_05_la_busqueda_completa.py}
\end{frame}

\marcoresultado{La búsqueda completa}{../figures/knapsack_05_la_busqueda_completa.png}{%
    El AG alcanza el valor exacto 50 con un peso de 20 con brecha cero en 4,060 evaluaciones}

% ================= radar_01_el_paisaje =================
\section{Un paisaje de cobertura}

\begin{frame}{Un paisaje de cobertura}
Cada bit decide si instalar un radar en un sitio candidato. Un radar
cubre un objetivo cuando su distancia euclidiana es a lo sumo el radio fijo.
\end{frame}

\marcoafirmacion{Un paisaje de cobertura}{%
    Nueve sitios definen 512 planes y juntos cubren los 36 objetivos}

% ================= radar_02_el_cromosoma =================
\section{El cromosoma}

\begin{frame}{El cromosoma}
El cromosoma
\end{frame}

\begin{frame}[fragile]{(1) Individuo}
\lstinputlisting[basicstyle=\ttfamily\fontsize{8pt}{9.6pt}\selectfont,firstline=39,lastline=49]{../src/radar_02_el_cromosoma.py}
\end{frame}

\begin{frame}[fragile]{(2) individuo\_aleatorio()}
\lstinputlisting[basicstyle=\ttfamily\fontsize{8pt}{9.6pt}\selectfont,firstline=52,lastline=59]{../src/radar_02_el_cromosoma.py}
\end{frame}

\begin{frame}[fragile]{(3) el censo}
\lstinputlisting[basicstyle=\ttfamily\fontsize{8pt}{9.6pt}\selectfont,firstline=62,lastline=66]{../src/radar_02_el_cromosoma.py}
\end{frame}

\marcoafirmacion{El cromosoma}{%
    21 de 500 planes aleatorios cubren todos los objetivos}

% ================= radar_03_la_penalizacion =================
\section{Una penalización lexicográfica}

\begin{frame}{Una penalización lexicográfica}
Una penalización lexicográfica
\end{frame}

\begin{frame}[fragile]{(1) objetivo()}
\lstinputlisting[basicstyle=\ttfamily\fontsize{8pt}{9.6pt}\selectfont,firstline=57,lastline=61]{../src/radar_03_la_penalizacion.py}
\end{frame}

\begin{frame}[fragile]{(2) la auditoría de clasificación}
\lstinputlisting[basicstyle=\ttfamily\fontsize{8pt}{9.6pt}\selectfont,firstline=64,lastline=70]{../src/radar_03_la_penalizacion.py}
\end{frame}

\marcoafirmacion{Una penalización lexicográfica}{%
    La penalización auditada clasifica a cada plan factible por encima de cada plan no factible}

% ================= radar_04_operadores =================
\section{Operadores binarios}

\begin{frame}{Operadores binarios}
Operadores binarios
\end{frame}

\begin{frame}[fragile]{(1) seleccion\_torneo()}
\lstinputlisting[basicstyle=\ttfamily\fontsize{8pt}{9.6pt}\selectfont,firstline=58,lastline=61]{../src/radar_04_operadores.py}
\end{frame}

\begin{frame}[fragile]{(2) cruza()}
\lstinputlisting[basicstyle=\ttfamily\fontsize{8pt}{9.6pt}\selectfont,firstline=64,lastline=68]{../src/radar_04_operadores.py}
\end{frame}

\begin{frame}[fragile]{(3) mutar()}
\lstinputlisting[basicstyle=\ttfamily\fontsize{8pt}{9.6pt}\selectfont,firstline=71,lastline=77]{../src/radar_04_operadores.py}
\end{frame}

\marcoafirmacion{Operadores binarios}{%
    Una generación reduce el objetivo medio de 86.0 a 42.6}

% ================= radar_05_la_busqueda_completa =================
\section{La búsqueda completa}

\begin{frame}{La búsqueda completa}
La búsqueda completa
\end{frame}

\begin{frame}[fragile]{(1) optimo\_exacto()}
\lstinputlisting[basicstyle=\ttfamily\fontsize{8pt}{9.6pt}\selectfont,firstline=76,lastline=80]{../src/radar_05_la_busqueda_completa.py}
\end{frame}

\begin{frame}[fragile]{(2) ejecutar()}
\lstinputlisting[basicstyle=\ttfamily\fontsize{8pt}{9.6pt}\selectfont,firstline=83,lastline=98]{../src/radar_05_la_busqueda_completa.py}
\end{frame}

\begin{frame}[fragile]{(3) el mapa de cobertura}
\lstinputlisting[basicstyle=\ttfamily\fontsize{8pt}{9.6pt}\selectfont,firstline=110,lastline=126]{../src/radar_05_la_busqueda_completa.py}
\end{frame}

\marcoresultado{La búsqueda completa}{../figures/radar_05_la_busqueda_completa.png}{%
    El AG y la enumeración exhaustiva encuentran cinco radares con cobertura completa}

% ================= schedule_01_los_requisitos =================
\section{Los requisitos}

\begin{frame}{Los requisitos}
Un bit asigna a un empleado a un turno. La instancia necesita 35 asignaciones
en 21 turnos, mientras que ocho empleados pueden cubrir como máximo cinco turnos cada uno.
\end{frame}

\marcoafirmacion{Los requisitos}{%
    168 bits codifican 35 asignaciones requeridas bajo la capacidad 40}

% ================= schedule_02_horarios_aleatorios =================
\section{Horarios aleatorios}

\begin{frame}{Horarios aleatorios}
Horarios aleatorios
\end{frame}

\begin{frame}[fragile]{(1) horario\_aleatorio()}
\lstinputlisting[basicstyle=\ttfamily\fontsize{8pt}{9.6pt}\selectfont,firstline=34,lastline=39]{../src/schedule_02_horarios_aleatorios.py}
\end{frame}

\begin{frame}[fragile]{(2) violaciones()}
\lstinputlisting[basicstyle=\ttfamily\fontsize{8pt}{9.6pt}\selectfont,firstline=42,lastline=56]{../src/schedule_02_horarios_aleatorios.py}
\end{frame}

\begin{frame}[fragile]{(3) el censo}
\lstinputlisting[basicstyle=\ttfamily\fontsize{8pt}{9.6pt}\selectfont,firstline=59,lastline=64]{../src/schedule_02_horarios_aleatorios.py}
\end{frame}

\marcoafirmacion{Horarios aleatorios}{%
    0 de 500 horarios aleatorios son factibles; el mejor todavía tiene 14 violaciones}

% ================= schedule_03_reparacion =================
\section{Reparar las restricciones}

\begin{frame}{Reparar las restricciones}
Reparar las restricciones
\end{frame}

\begin{frame}[fragile]{(1) costo\_preferencia()}
\lstinputlisting[basicstyle=\ttfamily\fontsize{8pt}{9.6pt}\selectfont,firstline=55,lastline=59]{../src/schedule_03_reparacion.py}
\end{frame}

\begin{frame}[fragile]{(2) reparar()}
\lstinputlisting[basicstyle=\ttfamily\fontsize{6pt}{7.2pt}\selectfont,firstline=62,lastline=94]{../src/schedule_03_reparacion.py}
\end{frame}

\begin{frame}[fragile]{(3) la comparación}
\lstinputlisting[basicstyle=\ttfamily\fontsize{8pt}{9.6pt}\selectfont,firstline=97,lastline=103]{../src/schedule_03_reparacion.py}
\end{frame}

\marcoafirmacion{Reparar las restricciones}{%
    La reparación hace factibles a 500/500 horarios; los costos reparados van desde 58 en adelante}

% ================= schedule_04_operadores =================
\section{Operadores bajo reparación}

\begin{frame}{Operadores bajo reparación}
Operadores bajo reparación
\end{frame}

\begin{frame}[fragile]{(1) seleccion\_torneo()}
\lstinputlisting[basicstyle=\ttfamily\fontsize{8pt}{9.6pt}\selectfont,firstline=81,lastline=84]{../src/schedule_04_operadores.py}
\end{frame}

\begin{frame}[fragile]{(2) cruza()}
\lstinputlisting[basicstyle=\ttfamily\fontsize{8pt}{9.6pt}\selectfont,firstline=87,lastline=92]{../src/schedule_04_operadores.py}
\end{frame}

\begin{frame}[fragile]{(3) mutar()}
\lstinputlisting[basicstyle=\ttfamily\fontsize{8pt}{9.6pt}\selectfont,firstline=95,lastline=106]{../src/schedule_04_operadores.py}
\end{frame}

\marcoafirmacion{Operadores bajo reparación}{%
    Todos los 100 hijos son factibles; el costo medio cae de 77.8 a 61.2}

% ================= schedule_05_la_busqueda_completa =================
\section{La búsqueda completa}

\begin{frame}{La búsqueda completa}
La búsqueda completa
\end{frame}

\begin{frame}[fragile]{(1) una\_generacion()}
\lstinputlisting[basicstyle=\ttfamily\fontsize{8pt}{9.6pt}\selectfont,firstline=102,lastline=109]{../src/schedule_05_la_busqueda_completa.py}
\end{frame}

\begin{frame}[fragile]{(2) ejecutar()}
\lstinputlisting[basicstyle=\ttfamily\fontsize{8pt}{9.6pt}\selectfont,firstline=112,lastline=123]{../src/schedule_05_la_busqueda_completa.py}
\end{frame}

\begin{frame}[fragile]{(3) el reporte base}
\lstinputlisting[basicstyle=\ttfamily\fontsize{8pt}{9.6pt}\selectfont,firstline=126,lastline=131]{../src/schedule_05_la_busqueda_completa.py}
\end{frame}

\marcoresultado{La búsqueda completa}{../figures/schedule_05_la_busqueda_completa.png}{%
    El AG con un costo de 41 vence al mejor de 5,100 horarios aleatorios reparados, con costo 52, con igual conteo de evaluaciones}

\section{Siguiente}

\begin{frame}{Siguiente}
El AG de la mochila da el valor exacto 50 a peso 20. El horario arranca en 0/500 factibles y aun así vence al aleatorio reparado a igual presupuesto.

\vspace{0.8em}
\textbf{Lección 10 --- Optimización combinatoria ordenada}

El orden y la unicidad ahora son parte de la corrección. El OX1 que en la 04 era un no-op silencioso es aquí el operador legal.
\end{frame}
\end{document}
